\documentclass[letterpaper, 11pt]{article}
\usepackage{latexsym}
\usepackage[empty]{fullpage}
\usepackage{titlesec}
\usepackage{marvosym}
\usepackage[usenames, dvipsnames]{color}
\usepackage{verbatim}
\usepackage{enumitem}
\usepackage[hidelinks]{hyperref}
\usepackage{fancyhdr}
\usepackage[T1]{fontenc}
\usepackage[english]{babel}
\usepackage{tabularx}
\usepackage{fontawesome5}
\usepackage{multicol}
\setlength{\multicolsep}{-3.0pt}
\setlength{\columnsep}{-1pt}
\input{glyphtounicode}

\hypersetup{
    colorlinks=true,
    linkcolor=black,
    filecolor=black,
    citecolor=black,
    urlcolor=RoyalBlue,
}

\pagestyle{fancy}
\fancyhf{}
\fancyfoot{}
\renewcommand{\headrulewidth}{0pt}
\renewcommand{\footrulewidth}{0pt}
\setlength{\footskip}{12pt}

% Adjust margins
\addtolength{\oddsidemargin}{-0.6in}
\addtolength{\evensidemargin}{-0.5in}
\addtolength{\textwidth}{1.19in}
\addtolength{\topmargin}{-.5in}
\addtolength{\textheight}{1.4in}

\urlstyle{same}

\raggedbottom
\raggedright
\setlength{\tabcolsep}{0in}

% Sections formatting
\titleformat{\section}{ \vspace{-4pt}\scshape\raggedright\large\bfseries }{}{0em}{}[
\color{black}
\titlerule
\vspace{-5pt}
]

\pdfgentounicode=1

%-------------------------
% Custom commands
\newcommand{\resumeItem}[1]{ \item\small{ {#1 \vspace{-2pt}} } }

\newcommand{\resumeSubheading}[4]{
\vspace{-1pt}
\item
\begin{tabular*}{1.0\textwidth}[t]{l@{\extracolsep{\fill}}r}
    \textbf{#1}       & {\small #2} \\
    {\small#3} & \textit{\small #4} \\
\end{tabular*}
\vspace{-10pt}
}

\newcommand{\resumeProjectHeading}[2]{
\vspace{-1pt}
\item
\begin{tabular*}{1.0\textwidth}[t]{l@{\extracolsep{\fill}}r}
    \textbf{#1}  & {\small #2} \\
\end{tabular*}
\vspace{-20pt}
}

\newcommand{\resumeSubItem}[1]{\resumeItem{#1}\vspace{-4pt}}

\renewcommand{\labelitemi}{$\vcenter{\hbox{\tiny$\bullet$}}$}
\renewcommand{\labelitemii}{$\vcenter{\hbox{\tiny$\bullet$}}$}

\newcommand{\resumeSubHeadingListStart}{\begin{itemize}[leftmargin=0.0in, label={}]}
\newcommand{\resumeSubHeadingListEnd}{\end{itemize}}
\newcommand{\resumeItemListStart}{\begin{itemize}[leftmargin=0.25in]}
\newcommand{\resumeItemListEnd}{\end{itemize}\vspace{-5pt}}

%-------------------------------------------
%%%%%%  RESUME STARTS HERE  %%%%%%%%%%%%%%%%%%%%%%%%%%%%

\begin{document}

%----------HEADING----------
\begin{center}
    {\Huge \scshape Gawtam Chithra Ramesh} \\
    \vspace{3pt}
    \small
    {
    \raisebox{-0.1\height}{\faPhone} +46-76541 8589 ~
    \href{mailto:gawtamcr3@gmail.com}{\raisebox{-0.2\height}{\faEnvelope}} {gawtamcr3@gmail.com} ~
    \href{https://linkedin.com/in/gawtamcr/}{\raisebox{-0.2\height}{\faLinkedin} \textcolor{black}{gawtamcr} } ~
    \href{https://github.com/gawtamcr}{\raisebox{-0.2\height}{\faGithub} \textcolor{black}{gawtamcr}}
    }
    \vspace{-8pt}
\end{center}

%-----------EDUCATION-----------
\section{Education}
    \resumeSubHeadingListStart
    \resumeSubheading{KTH Royal Institute of Technology}{Aug. 2024 -- Jul. 2026}{Master of Science in \href{https://www.kth.se/en/studies/master/systems-control-robotics/msc-systems-control-and-robotics-1.8733}{Systems, Controls, and Robotics}}{Stockholm, Sweden}
    \vspace{-0.4pt}
    \resumeItemListStart
        \resumeItem{\textbf{Coursework:} Robot Learning \& Embodied AI, Deep Learning Advanced, Image Analysis \& Computer Vision, Introduction to Robotics, Multi-Agent Systems.}
    \resumeItemListEnd

    \resumeSubheading{Indian Institute of Technology Madras}{Aug. 2019 -- Jul. 2024}{Dual Degree (B.Tech in \href{https://ed.iitm.ac.in/about.html}{Engineering Design} and M.Tech in \href{https://ed.iitm.ac.in/~robotics_lab/}{Robotics})}{Chennai, India}
    \vspace{-0.4pt}
    \resumeItemListStart
        \resumeItem{\textbf{Coursework:} Mechanics of Serial Robots, Field and Service Robotics, Control Systems, Fundamentals of Deep Learning, Data Structures and Algorithms.}
    \resumeItemListEnd
    \resumeSubHeadingListEnd

%-----------EXPERIENCE-----------
\section{Professional Experience}
\resumeSubHeadingListStart

    \resumeSubheading{Master's Thesis}{Jan 2026 -- Jun 2026}{\textbf{ABB Robotics}, Mentors: \href{https://www.linkedin.com/in/matthewwilliamlock/}{Matthew Lock} and \href{https://www.linkedin.com/in/jonathan-styrud-99385864/}{Jonathan Styrud}}{V\"aster\aa s, Sweden}
    \resumeItemListStart
        \resumeItem{Architecting a System 1 / System 2 framework for industrial manipulation: a generative flow-matching trajectory planner (System 1) steered at inference time by Signal Temporal Logic robustness gradients (System 2), replacing black-box policies with a verifiable, explainable intermediate layer.}
        \resumeItem{Designed a single generative model that can be steered by arbitrary STL specifications at inference time, eliminating retraining or per-task expert demonstrations and enabling flexible transfer across task variants.}
        \resumeItem{Validated the framework with full specification satisfaction on reach-avoid benchmarks; benchmarking head-to-head against state-of-the-art STL planners (DAG-STL, ZSTP) on Maze2D and AntMaze, targeting a top-venue robotics submission.}
        \resumeItem{Collaborating with ABB Robotics R\&D on transfer of the approach to long-horizon 6-DoF manipulation use cases.}
    \resumeItemListEnd

    \resumeSubheading{Device Research Intern, 3D Scene Graphs}{Jul 2025 -- Dec 2025}{\textbf{Ericsson Research}, Mentors: \href{https://se.linkedin.com/in/p\%C3\%BCren-g\%C3\%BCler-76457317}{P\"uren G\"uler} and \href{https://se.linkedin.com/in/hcaltenco}{Hector Caltenco}}{Lund, Sweden}
    \resumeItemListStart
        \resumeItem{Developed a robust 3D scene-graph pipeline by integrating instance segmentation and tracking into a SOTA framework (Hydra), reducing parameter sensitivity for reliable scene interpretation across diverse environments.}
        \resumeItem{Designed a synchronized RGB-D inpainting pipeline to remove privacy-sensitive objects from 3D scene graphs in real time, addressing trustworthy-AI requirements for human-centered deployment, and benchmarked SOTA inpainting algorithms balancing reconstruction fidelity with real-time performance.}
        \resumeItem{Lead author on submissions to \textbf{ICPR 2026} (under review) and \textbf{IMPROVE 2026} (accepted).}
    \resumeItemListEnd

    \resumeSubheading{Graduate Robotics Intern, Modular Collaborative Robots}{Dec 2022 -- Jul 2023}{\textbf{Systemantics India Pvt. Ltd.}, Mentor: \href{https://www.linkedin.com/in/jagannathraju/}{Jagannath Raju}}{Bangalore, India}
    \resumeItemListStart
        \resumeItem{Developed backward-chained Behavior Trees on a 6-DoF industrial manipulator for robust, long-horizon trajectory planning in collaborative-robot deployments.}
        \resumeItem{Designed a 6-DoF trajectory generation and control pipeline in C++, achieving real-time joint actuation via low-latency D-Bus and SocketCAN communication with motor controllers.}
    \resumeItemListEnd

    \resumeSubheading{Robotics Intern, Pallet Handling}{Jun 2021 -- Jul 2021}{\textbf{Yaskawa Pvt. Ltd.}, Mentor: Manju Tiwari}{Gurgaon, India}
    \resumeItemListStart
        \resumeItem{Programmed industrial 6-DoF manipulators for mixed-carton palletizing, a core machine-tending task, using MotoSim simulation and teach-pendant scripting.}
        \resumeItem{Built a C++ memory-optimization algorithm for streaming data between PC and robot controller, improving throughput in high-volume pick-and-place operations.}
    \resumeItemListEnd

\resumeSubHeadingListEnd

%-----------RESEARCH EXPERIENCE-----------
\section{Research Experience}
\resumeSubHeadingListStart

    \resumeSubheading{Graduate Research, Deformable Object Manipulation}{Jan 2025 -- Nov 2025}{\href{https://www.kth.se/is/rpl}{Robotics, Perception and Learning (RPL)}, KTH, Advisor: \href{https://www.kth.se/is/rpl}{Florian T. Pokorny}}{Stockholm, Sweden}
    \resumeItemListStart
        \resumeItem{Integrated taxonomy-guided vision-language models (Gemini, Qwen) to interpret deformable-object manipulation scenes and map structured outputs to motion primitives.}
        \resumeItem{Built a prompt and taxonomy framework that generates robot-parsable, formally-structured specifications from natural-language prompts; iteratively improved the framework through quantitative failure-mode analysis and integration of multimodal and temporal cues.}
        \resumeItem{Lead author on the resulting submission to the ROMADO workshop at \textbf{IROS 2025}.}
    \resumeItemListEnd
    
    \resumeSubheading{Dual Degree Project, Structure from Motion for Autonomous Ultrasound}{Jan 2024 -- May 2024}{\href{https://niravatgit.github.io/INSPIRELab/}{INSPIRE Lab}, Advisor: \href{https://niravatgit.github.io/patelnirav.com/index.html}{Nirav Patel}}{IIT Madras, India}
    \resumeItemListStart
        \resumeItem{Built a 3D perception system using Structure from Motion (COLMAP) on monocular RGB images for cost-efficient human torso reconstruction, enabling more accessible ultrasound diagnostics.}
        \resumeItem{Designed an autonomous scanning system integrating a 5-DoF robotic arm with single-camera RGB perception via ROS, MoveIt, and Gazebo simulation.}
        \resumeItem{Implemented an end-to-end perception-to-planning loop: 3D point-cloud generation, trajectory planning, and ultrasound-probe execution.}
    \resumeItemListEnd

\resumeSubHeadingListEnd

%-----------PROJECTS-----------
\section{Projects}
\resumeSubHeadingListStart

    \resumeProjectHeading{Robot Learning and Embodied AI $|$ \textnormal{\small DD2601, KTH}}{Sep 2025 -- Oct 2025}
    \resumeItemListStart
        \resumeItem{Built a 3D open-vocabulary semantic mapping pipeline using RGB-D data from ARKitScenes, integrating OwlV2 detection, SAM segmentation, and CLIP grounding with depth-fusion 2D-to-3D projection for queryable scene maps.}
        \resumeItem{Deployed a pre-trained Vision-Language-Action (VLA) model for robot manipulation, executing policy inference from visual inputs and natural-language commands; analyzed generalization and robustness limits on unseen tasks.}
    \resumeItemListEnd

    \resumeProjectHeading{Reinforcement Learning for Pacman Capture-the-Flag $|$ \textnormal{\small KTH}}{Mar 2025 -- May 2025}
    \resumeItemListStart
        \resumeItem{Designed a hybrid agent for Pacman Capture-the-Flag in Unity3D, merging a custom Behavior Tree for high-level strategy with role-specific PPO policies for specialized actions (attack, defend, escape) in a partially-observable, dynamic, multi-agent environment.}
        \resumeItem{Trained the RL policies via Unity ML-Agents using phased training, self-play, and curriculum learning, eliciting adaptive behaviour and on-the-fly role switching across competitive matches.}
    \resumeItemListEnd

    \resumeProjectHeading{Introduction to Robotics $|$ \textnormal{\small DD2410, KTH}}{Sep 2024 -- Oct 2024}
    \resumeItemListStart
        \resumeItem{Implemented an iterative inverse-kinematics solver from scratch for a 7-DoF KUKA arm in ROS.}
        \resumeItem{Integrated a Behavior Tree mission planner for a TIAGo mobile manipulator in ROS/Gazebo, enabling autonomous navigation, vision-based grasping, and recovery from kidnapped-robot scenarios.}
        \resumeItem{Engineered 2D occupancy-grid mapping in ROS from laser scans; applied RRT* for non-holonomic Dubins-car path planning.}
    \resumeItemListEnd

    \resumeProjectHeading{Multi-Agent RL for VR Shooter $|$ \textnormal{\small CFI, IIT Madras}}{May 2020 -- Aug 2020}
    \resumeItemListStart
        \resumeItem{Built multi-agent and cooperative-play PPO policies in Unity3D for a shooter strategy game with obstacles and procedurally generated mazes, achieving superhuman gameplay against human opponents.}
        \resumeItem{Integrated the trained agents into a Unity VR environment for direct virtual-reality interaction with the learned policies.}
    \resumeItemListEnd

\resumeSubHeadingListEnd

%-----------ACHIEVEMENTS-----------
\section{Achievements}
\begin{itemize}[leftmargin=0.15in]
    \setlength{\itemsep}{-0.3em}
    \item Awarded the \textbf{IIT Madras Young Research Fellowship} (top 30 of 800+ applicants) for outstanding research potential.
    \item Top-percentile performance in India's national entrance examinations: \textbf{97.64 percentile in JEE Advanced} and \textbf{99.10 percentile in JEE Mains} (out of 1.5 million applicants); \textbf{All-India Rank 52 of 50,000+} in UCEED (Design).
\end{itemize}

%-----------LEADERSHIP-----------
\section{Leadership \& Mentoring}
\resumeSubHeadingListStart

    \resumeProjectHeading{Head of Industrial \& Public Relations $|$ \textnormal{\small Department of Engineering Design, IIT Madras}}{Apr 2021 -- May 2022}
    \resumeItemListStart
        \resumeItem{Led a three-tier team running placement and internship processes for 1,200+ students; designed and executed the department's first hybrid recruitment framework, expanding alumni and industry outreach to broaden the opportunity pipeline.}
    \resumeItemListEnd

    \resumeProjectHeading{Strategist $|$ \textnormal{\small Computer Vision \& Intelligence Group, IIT Madras}}{Apr 2020 -- May 2021}
    \resumeItemListStart
        \resumeItem{Co-organized a Deep Learning summer school for 500+ students covering DL, RL, and computer vision; co-mentored student projects on image analysis and reinforcement learning for games.}
    \resumeItemListEnd

    \resumeProjectHeading{Graduate Teaching Assistant $|$ \textnormal{\small Field and Service Robotics (ID4060), IIT Madras}}{Aug 2023 -- Nov 2023}
    \resumeItemListStart
        \resumeItem{Designed and graded ROS and CoppeliaSim assignments for 40 students; ran small-group reviews on assignments and end-of-term projects across diverse student backgrounds.}
    \resumeItemListEnd

\resumeSubHeadingListEnd

%-----------SKILLS-----------
\section{Technical Skills}
\begin{itemize}[itemsep=-1pt, parsep=3pt, leftmargin=0.15in]
    \small{
        \item \textbf{Languages:} Python, C++, C, C\#  \hspace{1em} \textbf{Machine Learning:} PyTorch, JAX, OpenCV, wandb
        \item \textbf{Robotics \& Simulation:} ROS / ROS2, MoveIt, Gazebo, MuJoCo, PyBullet, Unity3D, Unity ML-Agents, CoppeliaSim, MATLAB
        \item \textbf{Tools:} Docker, Git, \LaTeX, CUDA, Slurm, VS Code
    }
\end{itemize}

\end{document}